
%\needspace{4\baselineskip}
\color{uniblue}{\subsection{Сигнализация замыканий на землю (СЗЗ)}}
\color{black}

\begin{enumerate}[label=\arabic{section}.\arabic{subsection}.\arabic*, labelsep=4pt, leftmargin=0pt, itemindent=57pt]

\item
СЗЗ предназначена для определения замыканий на землю в сетях с компенсированной или изолированной нейтралью. В устройстве сигнализация замыканий на землю представлена функцией «СЗЗ».

\item
Функция фиксирует наличие замыкания на землю по повышению напряжения $3U_0$ выше заданной уставки и, через выдержку времени, определяемую уставкой таймера Т1, выдает сигнализацию замыкания на землю.

\item
При повышении напряжения $3U_0$ выше уставки и отсутствии блокирующих сигналов происходит пуск защиты. По истечении выдержки времени Т1, в течение которой величина напряжения $3U_0$ постоянно была выше уставки срабатывания, происходит срабатывание защиты. Прохождение сигнала срабатывания СЗЗ блокируется при появлении сигнала неисправности цепей напряжения БНН, если введена уставка (SGF2 «Реж\_БНН» в положении «С контролем»).

\item
Ввод функции в работу на этапе параметрирования устройства осуществляется программным переключателем SGF1 «Ввод\_функции». Информация о введенной функции отображается выходным сигналом «СЗЗ: Ввод».

\item
Функция может быть выведена из работы оперативно путем активации сигнала «ОВ СЗЗ». Выведенное состояние функции характеризуется активным состоянием сигнала «СЗЗ: Оперативный вывод».

\item
Алгоритм работы функции выполнен в соответствии с рисунком \ref{szz:img1}. В таблице \ref{szz:tbl1} приведены параметры, необходимые для настройки функции.

\vspace{3mm}
\begin{figure}[H]
\centering
\includegraphics[width=0.8\textwidth,height=0.8\textheight,keepaspectratio]{\fbpath/203. СЗЗ/_latex/img/szz.pdf}
\captionof{figure}{Функциональная схема логики СЗЗ}\label{szz:img1}
\end{figure}


\end{enumerate}
